\section{Discussion}
\label{sec:discussion}

\paragraph{Is mining bug reports better than generating SQL?} Not on this
target. The generator reached more defects at the same budget
(Observation~\ref{obs:gen}), and the templates are blocked from defect~4 by their
ASCII string fills. Mining did two things the generator did not. It found
defects~1 and~3 first in time, before the generator's vocabulary contained
\sql{lpad} or \sql{repeat}; we added those functions to the generator only after
the defects were filed. And it pointed the work at pad functions in the first
place. The fair reading is that each source covers a weakness of the other, which
is why drawing from both reached a defect soonest
(Observation~\ref{obs:both}).

\paragraph{Are four defects in one family really four?} They are four ways for a
user to end the process, each from a different statement, so we count them as
four. They are not four independent mistakes: defects 1 and 2 share one
function body, and defects 1 to 3 share one cause and one patch. Upstream records
them as two issues and two pull requests. We report both counts and say which is
which, so that neither inflates the other.

\paragraph{Would more fuzzing time have found them?} No.
Observation~\ref{obs:blind} is a property of the fuzzer's source, not of its
budget. A generator cannot emit a function outside its vocabulary or a value
outside its domain however long it runs. This is the main practical point of the
paper: the incumbent fuzzer's run time tells nothing about how close it is to
these defects.

\paragraph{Why did fixing three defects reveal a fourth?} Before the first patch,
the length check ran after the allocation, so its unit did not matter: the
process died first. Once the check ran first, its unit became the thing that
decided the outcome. A patch is a good place to look for the next defect, since
it marks code where the authors' first understanding was wrong. Our own first
understanding of defect~4 was wrong too, and a reviewer corrected it.

\paragraph{Does this work on other engines?} We do not know. The method needs a
second engine with many closed, labelled reports that carry reproductions, and
an engine under test that shares enough of its dialect for some reports to run.
Only 16 of 60 reports that we classified by hand looked transferable to another
engine. Testing other pairs of engines, and filling string holes with
multi-byte values, we leave for future work.
